% Korean edition of figs/signalr.tex.
% Time-space diagram in the style of Kurose & Ross Fig 2.7/3.39: two timelines, time runs down, sloped arrows,
% message names along the arrows, events beside the timelines. The time axis is broken twice (zigzags).
\fig[wide]{signalr}{협상과 핸드셰이크를 마친 뒤 서버는 피드 요청에만 답한다}[클라이언트와 실시간 피드 허브 \code{/hubs/main} 사이의 메시지. 시간은 아래로 흐른다. 연결 설정은 밀리초 단위, 그다음은 3\,s 폴링 주기]{%
\begin{tikzpicture}[x=1mm,y=1mm,
  f/.style={font=\fontsize{7}{8.4}\selectfont},
  host/.style={font=\fontsize{7.5}{9}\selectfont\bfseries,text=navy,anchor=south,align=center},
  tl/.style={greyM,line width=0.6pt},
  msg/.style={-{Stealth[length=1.6mm]},line width=0.7pt,navy},
  rep/.style={-{Stealth[length=1.6mm]},line width=0.7pt,greyD},
  ml/.style={font=\fontsize{7.5}{9}\selectfont,sloped,above,inner sep=0.7pt,text=navy},
  rl/.style={font=\fontsize{7.5}{9}\selectfont,sloped,below,inner sep=0.7pt,text=greyD},
  ev/.style={f,anchor=east,align=right,text=greyD},
  brk/.style={line width=0.6pt,greyM,fill=white},
  stp/.style={f,anchor=west,text width=48mm,align=left,text=greyD,execute at begin node={\hyphenpenalty=10000\exhyphenpenalty=10000\relax}},
  brace/.style={decorate,decoration={brace,amplitude=1mm},line width=0.5pt,navy}]
\def\C{0}\def\S{92}
% ---------------------------------------------------------------- hosts and timelines
\node[host] at (\C,1) {클라이언트\\(브라우저 또는 래퍼)};
\node[host] at (\S,1) {서버\\(SignalR 허브)};
\draw[tl] (\C,0) -- (\C,-106); \draw[tl] (\S,0) -- (\S,-106);
\draw[tl,-{Stealth[length=1.4mm]}] (\C,-106) -- (\C,-109); \draw[tl,-{Stealth[length=1.4mm]}] (\S,-106) -- (\S,-109);
\node[f,text=greyD,anchor=north] at (\C,-109.3) {시간};\node[f,text=greyD,anchor=north] at (\S,-109.3) {시간};
% broken-axis marks: two zigzags across both timelines
\newcommand{\sigzig}[1]{\foreach \x in {\C,\S}{\fill[white] (\x-1.6,#1+0.9) rectangle (\x+1.6,#1-0.9);
  \draw[greyM,line width=0.5pt] (\x-1.6,#1+1.4) -- (\x+1.6,#1+0.4); \draw[greyM,line width=0.5pt] (\x-1.6,#1-0.4) -- (\x+1.6,#1-1.4);}}
% ---------------------------------------------------------------- set-up (milliseconds)
\draw[msg] (\C,-6) -- node[ml]{\code{POST /hubs/main/negotiate?negotiateVersion=1}} (\S,-10);
\draw[rep] (\S,-11) -- node[rl]{200: \code{connectionId}, \code{connectionToken}, 전송 방식 3개} (\C,-15);
\draw[msg] (\C,-23) -- node[ml]{\code{GET /hubs/main?id=...}, \code{Upgrade: websocket}} (\S,-27);
\draw[rep] (\S,-28) -- node[rl]{\code{101 Switching Protocols}} (\C,-32);
\draw[msg] (\C,-40) -- node[ml]{\code{{"protocol":"json","version":1}} 0x1E} (\S,-44);
\draw[rep] (\S,-45) -- node[rl]{\code{{}} 0x1E} (\C,-49);
% events beside the client timeline
\node[ev] at (\C-2,-8) {연결 1, HTTP/2,\\REST용으로 이미 열림};
\node[ev] at (\C-2,-25) {연결 2: 새 TLS\\연결, HTTP/1.1};
\node[ev] at (\C-2,-42) {핸드셰이크};
% ---------------------------------------------------------------- poll cycle
\draw[msg] (\C,-56) -- node[ml]{\code{RequestPortfolio} \code{{intv:3, auto:true, lastT:0, d, pfId}}} (\S,-60);
\draw[rep] (\S,-61) -- node[rl]{\code{ReceivePortfolio}: 스냅샷 (\code{bs}, \code{ords}, \code{pos}, \code{msecs})} (\C,-65);
\draw[brace] (\C-2,-65) -- (\C-2,-56) node[midway,ev,xshift=-1.2mm] {응답\\20\textasciitilde 84\ms};
% 3 s wait with a break
\fill[navyM] (\C-0.6,-65.5) rectangle (\C+0.6,-76.5);
\sigzig{-71}
\draw[brace] (\C-2,-76.5) -- (\C-2,-65.5) node[midway,ev,xshift=-1.2mm] {3\,s 대기\\(축척 아님)};
\draw[msg] (\C,-77) -- node[ml]{\code{RequestPortfolio} \code{{..., lastT = t of previous reply}}} (\S,-81);
\draw[rep] (\S,-82) -- node[rl]{\code{ReceivePortfolio}: 델타 (\code{pos_diff}, \code{ords_diff}, \code{msecs_diff})} (\C,-86);
\fill[navyM] (\C-0.6,-86.5) rectangle (\C+0.6,-93);
\sigzig{-90}
\node[ev] at (\C-2,-90) {3\,s 대기 후 반복};
% ping: both directions, drawn level (schematic)
\draw[{Stealth[length=1.6mm]}-{Stealth[length=1.6mm]},line width=0.6pt,greyD,dashed] (\C,-100) -- node[f,above,inner sep=0.8pt,text=greyD]{유형 6 핑, 양쪽 모두 보냄. 서버는 15\,s마다} (\S,-100);
\node[ev] at (\C-2,-100) {유휴 연결을\\열어 둠};
% ---------------------------------------------------------------- step brackets on the right
\foreach \a/\b/\n/\t in {-5/-16/1/{\textbf{협상}: REST 연결에서 Bearer 토큰을 실은 HTTP/2 POST. 응답은 WebSocket, Server-Sent Events, 롱 폴링을 제시},
  -22/-33/2/{\textbf{연결}: 서버가 HTTP/2 위의 WebSocket[1]을 지원하지 않아 업그레이드가 두 번째 TLS 연결을 연다. 래퍼는 토큰을 헤더로 보내고, 브라우저는 URL에 넣는다},
  -39/-50/3/{\textbf{핸드셰이크}: 양쪽이 JSON 메시지에 합의하며, 메시지는 각각 0x1E 바이트로 끝난다},
  -55/-95/4/{\textbf{폴링}: 서버는 요청에 대한 답으로만 포트폴리오를 보낸다. 클라이언트는 응답마다 3\,s 뒤 다시 요청한다. \code{lastT}가 직전 응답의 \code{t}와 같으면 델타가 온다(\figref{lastt})}} {
  \draw[greyM,line width=0.5pt] (\S+3,\a) -- (\S+4,\a) -- (\S+4,\b) -- (\S+3,\b);
  \node[stp] at (\S+5.5,{(\a+\b)/2}) {\n.\ \t};
}
\end{tikzpicture}
}{남색 화살표는 클라이언트가, 회색 화살표는 서버가 보낸다. 시간은 아래로 흐른다. 연결 설정 부분은 측정한 지속 시간 없이 밀리초 단위로 그렸고, 지그재그 두 개는 줄여 그린 3\,s 대기를 뜻한다. 응답 시간 20\textasciitilde 84\ms{}는 새벽에 \code{auto}를 true와 false로 각각 두고 쟀다. 양쪽이 서로 독립적으로 핑을 보내므로 핑 화살표는 수평으로 그렸다. 유형 1은 상대편 메서드 호출, 유형 6은 핑, 유형 7은 종료다. 2026년 10월 28일 브라우저 기록에서 클라이언트는 직전 응답 3\,s 뒤에 다음 요청을 보냈다.}{브라우저 개발자 도구, Python 구현, DART Quant Team}
